\clearpage
\VerdierPhysicalPageAnchor{0137}%
\begingroup
\small
\setlength{\emergencystretch}{1.5em}%
\setlength{\abovedisplayskip}{1pt plus 1pt minus 1pt}%
\setlength{\belowdisplayskip}{1pt plus 1pt minus 1pt}%

\VerdierSemanticAnchor{ch2-2-2-2-sm2-completion-diagram}
\noindent into a commutative diagram:

\begin{center}
  \begin{tikzpicture}[>=stealth,line width=.38pt,x=1.30cm,y=1.45cm,
    scale=.92,transform shape]
    \node (Zpp) at (0,2.20) {$Z''$};
    \node (Xp)  at (-1.25,1.10) {$X'$};
    \node (Yp)  at ( 1.25,1.10) {$Y'$};
    \node (Y)   at (0,0) {$Y$};
    \draw[->] (Zpp) -- node[above left=-1pt] {$t'\in S$} (Xp);
    \draw[->] (Zpp) -- node[above right=-1pt] {$m'$} (Yp);
    \draw[->] (Xp)  -- node[below left=-1pt] {$m$} (Y);
    \draw[->] (Yp)  -- node[below right=-1pt] {$t\in S$} (Y);
    \node at (2.50,-.05) {$.$};
  \end{tikzpicture}
\end{center}

\VerdierSemanticAnchor{ch2-2-2-2-composite-roof-d3}
\noindent We then obtain a diagram
$D_3=X\xleftarrow{st'}Z''\xrightarrow{pm'}Z$, where the composite $st'$
belongs to $S$ (SM1).
\VerdierSemanticAnchor{ch2-2-2-2-composite-class-independence}%
One checks that the class $\overline D_3$ of $D_3$ in
$\operatorname{Hom}_{\mathcal C(S^{-1})}(X,Z)$ depends only on the classes
$\overline D_1$ and $\overline D_2$.
\VerdierSemanticAnchor{ch2-2-2-2-composition-conclusion}%
The morphism $\overline D_3$ is the
composite in $\mathcal C(S^{-1})$ of the morphisms $\overline D_1$ and
$\overline D_2$.

\VerdierSemanticAnchor{ch2-2-2-3}
\smallskip
\noindent\textbf{Remark 2.2.3.}
\VerdierSemanticAnchor{ch2-2-2-3-a}%
a) It follows from the fourth assertion of theorem (2.2.2) that the category
$\mathcal C(S^{-1})$, equipped with the functor
$Q\colon\mathcal C\longrightarrow\mathcal C(S^{-1})$, is a solution of a
universal problem:
\VerdierSemanticAnchor{ch2-2-2-3-a-universal-problem}%
the functor $Q$ carries the morphisms of $S$ to isomorphisms, and every
functor from $\mathcal C$ to a category $\mathcal A$ that carries the
morphisms of $S$ to isomorphisms factors uniquely through $Q$.
\VerdierSemanticAnchor{ch2-2-2-3-a-unique-isomorphism}%
The category $\mathcal C(S^{-1})$, equipped with the functor $Q$, is
therefore determined up to unique isomorphism by this property.

\VerdierSemanticAnchor{ch2-2-2-3-b}
\par\noindent b) Let $\mathcal C$ be a category and $S$ a set of morphisms
of $\mathcal C$ that need not be a multiplicative system of $\mathcal C$.
\VerdierSemanticAnchor{ch2-2-2-3-b-general-localization}%
There always exist a category $\mathcal C(S^{-1})$ and a functor
$Q\colon\mathcal C\longrightarrow\mathcal C(S^{-1})$ that carries the
morphisms of $S$ to isomorphisms, such that every functor from $\mathcal C$
to a category $\mathcal A$ carrying the morphisms of $S$ to isomorphisms
factors uniquely through $Q$.
\VerdierSemanticAnchor{ch2-2-2-3-b-multiplicative-comparison}%
When $S$ is a multiplicative system of $\mathcal C$, theorem (2.2.2) gives a
simple description of the category $\mathcal C(S^{-1})$; it also gives
additional properties of the functor $Q$ (assertions b) and c)) that do not
hold in general.

\VerdierSemanticAnchor{ch2-2-2-4}
\smallskip
\noindent\textbf{Corollary 2.2.4.}\quad
\begingroup\itshape
Let $\mathcal C$ be a category and $S$ a multiplicative system of
$\mathcal C$.

\VerdierSemanticAnchor{ch2-2-2-4-a}
\par\noindent a) The set $S$ is a multiplicative system of
$\mathcal C^\circ$, the category opposite to $\mathcal C$.
\VerdierSemanticAnchor{ch2-2-2-4-a-opposite-system}%
There is a canonical isomorphism of categories:

\VerdierSemanticAnchor{ch2-2-2-4-a-canonical-isomorphism}
\[
  \bigl(\mathcal C^\circ(S^{-1})\bigr)^\circ
  \longrightarrow \mathcal C(S^{-1})\,.
\]

\VerdierSemanticAnchor{ch2-2-2-4-a-hom-colimit-isomorphism}
\noindent This isomorphism induces in particular isomorphisms:
\[
  \operatorname{Hom}_{\mathcal C(S^{-1})}(X,Y)
  \longrightarrow
  \varinjlim_{Y\backslash S}\operatorname{Hom}_{\mathcal C}(X,\mathord\cdot)\,.
\]
\endgroup

\VerdierSemanticAnchor{ch2-2-2-4-p138-boundary}
\endgroup
